% !TeX root = ../../Book1.tex
% 纲要标题：### E.3 多项式的实根与实代数
% 纲要位置：计划纲要.md，第 930 行。
% BEGIN APPENDIX OUTLINE SYNC
% #### E.3.1 中间值性质与 Sturm 根计数
% #### E.3.2 传递原理、量词消去与半代数集合
% ---
% END APPENDIX OUTLINE SYNC

\section{多项式的实根与实代数}
\label{b1:sec:E03}

实根计数可以只用多项式除法与符号比较完成。本节先从实闭域上的因式分解推出中间值性质，再给出 Sturm 序列的计数公式。即使域中有大于所有整数的元素，证明也仍然适用。

% 纲要内容（仅作写作计划，尚非教材正文）：
%
% #### E.3.1 中间值性质与 Sturm 根计数
% #### E.3.2 传递原理、量词消去与半代数集合
% ---
%

\subsection{中间值性质与 Sturm 根计数}
\label{b1:subsec:E03:01}

\begin{theorem}[多项式中间值性质]\label{b1:thm:real-closed-ivt}
设 \(R\) 实闭，\(a<b\) 属于 \(R\)，且 \(f\in R[X]\) 满足 \(f(a)f(b)<0\)。则存在 \(c\in R\)，使 \(a<c<b\) 且 \(f(c)=0\)。
\end{theorem}

\begin{proof}
由\cref{b1:cor:real-polynomial-factorization}，将 \(f\) 分成常数、一次因子和处处为正的二次因子。若没有实根位于 \((a,b)\)，每个一次因子在两端的符号相同，二次因子在两端均为正，故乘积 \(f(a)f(b)>0\)，矛盾。
\end{proof}

例如，在实闭域中 \(f=X^3-X-1\) 满足 \(f(1)=-1\)、\(f(2)=5\)，所以 \((1,2)\) 中至少有一个根。仅看端点异号还不知道有几个根；如果两端同号，也不能断言无根，例如 \((X-1)(X-2)\) 在 \(0\) 和 \(3\) 处都为正，却在中间有两个根。

\ProseParagraphBreak
因式分解还说明：非零多项式的符号在不含根的区间上固定。根计数因此可以转化为研究越过各个根时符号怎样改变。Sturm 方法把 \(f\)、\(f'\) 和一串负余数放在一起：\(f'\) 帮助识别 \(f\) 的简单根，而余数前的负号使其他多项式的零点不会改变总变号数。以下把这两个作用分别证明。

\ProseParagraphBreak
先设 \(f\in R[X]\) 为无重根的非常数多项式。它与 \(f'\) 互素：若某个因子同时整除两者，就会给出重根。令 \(f_0=f\)、\(f_1=f'\)，将 \(f_{j-1}\) 除以 \(f_j\)，把所得非零余数的相反数记为 \(f_{j+1}\)，即
\begin{equation}\label{b1:eq:sturm-remainders}
f_{j-1}=q_jf_j-f_{j+1},\qquad \deg f_{j+1}<\deg f_j,
\end{equation}
如果余数为零便停止，记最后一个非零多项式为 \(f_m\)。余数次数严格下降保证过程终止；Euclid 算法说明最后一项与 \(\gcd(f,f')\) 相差一个非零常数因子，所以互素性保证 \(f_m\) 是非零常数。所得序列称为 \(f\) 的\term[Sturm-xulie]{Sturm 序列}[Sturm sequence]。对 \(x\in R\)，把 \(f_0(x),\ldots,f_m(x)\) 中的零项删去，记剩余符号列的变号次数为 \(V(x)\)：每当相邻两个非零值的符号相反，就计一次。例如符号列 \((+,0,-,-,+)\) 先删成 \((+,-,-,+)\)，共有两次变号，而不是把零另当作第三种符号。
\symidx{\(V(x)\)：Sturm 序列在一点的变号次数}

\begin{theorem}[Sturm]\label{b1:thm:sturm-count}
在上述条件下，若 \(a<b\) 且 \(f(a)f(b)\neq0\)，则 \(f\) 在 \((a,b)\) 中的根数为 \(V(a)-V(b)\)，每个根只计一次。
\end{theorem}

\begin{proof}
相邻两个 \(f_j\) 不可能有公共根。若 \(f_j(c)=f_{j+1}(c)=0\)，把递推式逐项向后使用，会迫使 \(f_{j+2}(c),\ldots,f_m(c)\) 都为零，最终与 \(f_m\) 是非零常数矛盾。

先考察一个中间项的根，即 \(f_j(c)=0\)、\(1\leq j<m\)。递推式给出 \(f_{j-1}(c)=-f_{j+1}(c)\neq0\)。取 \(c\) 两侧足够近的位置，使两个邻项都不经过自己的根，则其符号仍相反。假如邻项分别是正、负，中间项可能使三项成为 \((+,+,-)\) 或 \((+,-,-)\)，两者都只贡献一次变号；邻项为负、正时结论相同。在 \(c\) 处删去零项后，两邻项也恰贡献一次。因此经过中间项的零点不改变总变号数。若多个中间项在 \(c\) 处为零，它们不相邻；各自涉及的相邻符号对不重复，所以上述计算可以同时使用。

再考察 \(f\) 的根 \(c\)。无重根性给出 \(f'(c)\neq0\)。写成 \(f(X)=(X-c)g(X)\)，求导并代入 \(c\) 得 \(g(c)=f'(c)\neq0\)。在不跨越 \(g\) 和 \(f'\) 其他根的小区间里，两者符号不变且相同。\(c\) 左侧的 \(X-c\) 为负，所以 \(f\) 与 \(f'\) 异号；右侧的 \(X-c\) 为正，所以两者同号。因此从左到右经过 \(c\)，最前一对的变号次数由一变成零，恰减少一；后面各项按前段不改变总变号数。

这些局部比较并不需要分析学中的连续性定理。每个非零多项式只有有限个根，将全部 \(f_j\) 的根按序排列，相邻根之间的各项符号由\cref{b1:cor:real-polynomial-factorization}保持不变。在需要选择比较点时，取两个相邻根的中点即可；在端点外侧也可以取半个正距离。由此从 \(a\) 到 \(b\) 逐段比较，变号数只在 \(f\) 的根处各减少一。端点若有中间项为零，前段的删零规则保证 \(V(a)\)、\(V(b)\) 与邻近区间上的计数相同。总减少量因而恰为 \((a,b)\) 内 \(f\) 的根数。
\end{proof}

若原多项式有重根，先以 \(h=f/\gcd(f,f')\) 代替 \(f\)。实闭域特征为零；对每个不可约因子比较导数中的重数，可知 \(h\) 恰保留各因子一次，因此与 \(f\) 有相同的互异根，且无重根。上述算法于是计算互异实根数；若要计算重数，另保留原分解中的指数。

\ProseParagraphBreak
回到 \(f=X^3-X-1\)。先取 \(f_1=3X^2-1\)，逐次相除得
\[
\begin{aligned}
X^3-X-1&=\frac X3(3X^2-1)-\left(\frac23X+1\right),\\
3X^2-1&=\left(\frac92X-\frac{27}4\right)
\left(\frac23X+1\right)+\frac{23}4.
\end{aligned}
\]
第二式的余数是 \(23/4\)，其相反数才是最后的 Sturm 项。因此得到
\[
f_0=X^3-X-1,\quad f_1=3X^2-1,\quad
f_2=\tfrac23X+1,\quad f_3=-\tfrac{23}{4}.
\]
在 \(1\) 处四个值为 \(-1,2,5/3,-23/4\)，符号列是 \((-,+,+,-)\)，所以 \(V(1)=2\)。在 \(2\) 处四个值为 \(5,11,7/3,-23/4\)，符号列是 \((+,+,+,-)\)，所以 \(V(2)=1\)。于是 \((1,2)\) 内恰有一个实根，这比中间值性质提供的“至少一个”更精确。

\ProseParagraphBreak
同样，在 \(-2\) 处符号列为 \((-,+,-,-)\)，所以 \(V(-2)=2\)，\((-2,2)\) 中也只包含这一个根。当 \(x\geq2\) 时，\(x(x^2-1)-1\geq5\)；当 \(x\leq-2\) 时，该式不大于 \(-7\)。故外侧没有根，已经数出了全部实根。以上仅用有理数运算和符号比较，在每个实闭域中都成立。一般使用 Sturm 方法时，若工作域中的系数能精确运算并能判定符号，便可依此执行算法；抽象域的定义本身不保证存在可用的机器表示。

\subsection{传递原理、量词消去与半代数集合}
\label{b1:subsec:E03:02}

一个\term[bandaishu-jihe]{半代数集合}[semialgebraic set]是实闭域 \(R\) 的 \(R^n\) 中，由有限多个多项式等式与严格不等式通过有限次并、交、补得到的集合。系数所在的域须随问题指定；例如圆盘可写为 \(x^2+y^2\leq1\)，其中非严格不等式可拆成严格不等式与等式。

\ProseParagraphBreak
所谓投影，是只保留一些坐标而忘去其余坐标。若 \(S\subseteq R^{n+1}\) 由条件 \(\Phi(x_1,\ldots,x_n,y)\) 描述，那么其前 \(n\) 个坐标的投影由“存在 \(y\) 使 \(\Phi\) 成立”描述。消去这个存在量词，就是把条件改写成只涉及 \(x_1,\ldots,x_n\) 的有限多个等式与不等式。例如
\[
\exists y\ (y^2=x)\quad\Longleftrightarrow\quad x\geq0
\]
在实闭域中成立：平方非负给出必要性，正元素有平方根给出充分性。对一般有限多项式条件也能消去量词，是实闭域的\term[liangci-xiaoqu]{量词消去}[quantifier elimination]定理；这里仅预告其内容，不证明一般消去算法，也不把它当作上述 Sturm 证明的前提。

\ProseParagraphBreak
这里的一阶公式由域元素之间的等式、不等式，经过“且、或、非”以及对域元素的“存在”“任意”量词组成。被量词约束的是域元素，而不是任意子集；未被量词约束的变量称为自由变量。没有自由变量的公式是句子，例如“每个正元素都有平方根”可以写成
\[
\forall x\,\bigl(x>0\ \Longrightarrow\ \exists y\ (y^2=x)\bigr).
\]
若公式还点名使用某个给定域元素，例如 \(x^2>a\) 中事先指定的 \(a\)，就称使用了参数 \(a\)。这些说明只为读懂下一项预告，不要求先学习完整的逻辑理论。

\ProseParagraphBreak
相应的\term[chuan-di-yuanli]{传递原理}[transfer principle]说：用 \(0,1,+,-,\cdot,<\) 表述、没有额外参数的一阶句子，在所有实闭域中真值相同；若实闭域 \(R\subseteq S\)，则带 \(R\) 中参数的一阶公式，在取自 \(R\) 的自由变量值上也有相同真值。这同样作为进一步理论的预告，不在此证明。它不表示任意关于实数的分析陈述都能转移。例如“每个非空有上界的子集都有上确界”量化的是任意子集，不是这里的一阶域语言，而\secref{b1:sec:E02}已给出该性质失效的实闭域。

\ProseParagraphBreak
半代数集合的投影、量词消去和正性证书将转第三卷附录 G；相应的 Tarski--Seidenberg 原理、投影定理与正性定理可参阅 Bochnak、Coste 与 Roy 的《Real Algebraic Geometry》\cite[Sections 1.4, 2.2, 4.4]{BochnakCosteRoy1998RealAG}。这里不展开模型论、赋值理论或二次型分类。读者若只需要核验一元多项式的实根数，前一小节的除法表、端点符号和\cref{b1:thm:sturm-count}已经组成独立可核验的计算过程。
